
\chapter{Introdução} \label{cap:intro}

Veículos modernos são compostos por uma variedade de sistemas eletrônicos responsáveis pelo monitoramento e pelo controle de diversos subsistemas do automóvel \cite{sonko2024embedded}. Nesse contexto, destaca-se o sistema On-Board Diagnostics II (OBD-II), um padrão de diagnóstico embarcado que permite o monitoramento de parâmetros do veículo e a identificação de falhas nos seus componentes eletrônicos e mecânicos \cite{csselectronics_obd2_intro}. 

O OBD-II foi inicialmente desenvolvido com o objetivo de facilitar a detecção de falhas relacionadas às emissões de poluentes, permitindo um diagnóstico mais rápido e preciso \cite{sonko2024embedded}.  Ao longo do tempo, tornou-se um componente essencial na arquitetura de veículos modernos, sendo utilizado para monitoramento de sensores e leitura de códigos de erros \cite{csselectronics_obd2_intro}. No Brasil, a adoção deste sistema tornou-se obrigatória para todos os veículos leves a combustão produzidos ou importados a partir de 2011 por meio da Resolução CONAMA nº 354/2004 \cite{conama354_2004}, o que evidencia sua prevalência e sua importância.

Com a popularização do OBD-II, surgiram dispositivos conhecidos como \textit{dongles} OBD-II, que se conectam à porta de diagnóstico do veículo. Esses sistemas podem coletar dados do veículo em tempo real e transmiti-los para outros sistemas, normalmente através de comunicações sem fio como o Bluetooth ou o Wi-fi \cite{wen2020plugnpwned}. Frequentemente, estes \textit{dongles} são utilizados em conjunto com aplicativos móveis, conhecidos como \textit{companion apps}, que permitem analisar, visualizar e processar os dados coletados, ou até mesmo enviar mensagens ao barramento CAN \cite{companionapps2024security}.

Esta combinação de \textit{dongles} e \textit{companion apps} abriu espaço para o desenvolvimento de diversas aplicações no contexto automotivo. Com a coleta de dados através de \textit{dongles} conectados à porta OBD-II, constróem-se sistemas de gerenciamento de frotas, estabelecendo-se um monitoramento operacional e logístico  \cite{malekian2017fleet}; sistemas de identificação de motorista, que classificam os motoristas em perfis de condução baseados em sua “agressividade” \cite{barreto2018perfis}, o que por sua vez abre espaço para a oferta de seguros com prêmio personalizado \cite{michailidis2025mlobd}; e técnicas de manutenção preventiva  \cite{errezgouny2025predictive}.

A maioria destas aplicações pressupõe que o \textit{dongle} permaneça conectado ininterruptamente à porta OBD-II, de modo a coletar os dados continuamente. Esta necessidade contrasta com evidências de que diversos \textit{dongles} OBD-II disponíveis comercialmente apresentam vulnerabilidades de segurança \cite{wen2020plugnpwned}. Além disso, os próprios \textit{companion apps} associados a esses dispositivos também apresentam vulnerabilidades, possibilitando a realização de ataques mesmo por indivíduos que não possuam conhecimento aprofundado sobre o funcionamento interno dos sistemas envolvidos \cite{companionapps2024security}.

%%%%%%%%%%%%%%%%%%%%%%%%%%
\section{Objetivos}

O presente trabalho propõe o projeto e a implementação de um \textit{dongle} OBD-II e de
um \textit{companion app} com foco em segurança. O objetivo é desenvolver uma solução
que mantenha a funcionalidade esperada (leitura de parâmetros e DTCs) ao mesmo
tempo em que incorpore mecanismos robustos de autenticação, criptografia, proteção de
firmware e controle de acesso ao barramento CAN, mitigando as vulnerabilidades
catalogadas pelo \textit{Donglescope} \cite{wen2020plugnpwned}, e habilitando a utilização segura
do OBD-II para diferentes aplicações, como gerenciamento de frotas \cite{malekian2017fleet}, identificação do perfil do motorista \cite{barreto2018perfis}, personalização de seguro \cite{michailidis2025mlobd}, e técnicas de manutenção preventiva  \cite{errezgouny2025predictive}.

Almeja-se um protótipo funcional de um \textit{dongle} OBD-II e de um \textit{companion app} projetados com foco em segurança. A solução deverá ser capaz de realizar a leitura de parâmetros padronizados do veículo obtidos por meio da interface OBD-II, tais como códigos de erro (DTCs) gerados pelo sistema de autodiagnóstico do carro e dados de sensores e variáveis operacionais do motor, incluindo rotação do motor, velocidade do veículo, temperatura do líquido de arrefecimento e posição do acelerador. Esses dados serão acessados por meio de \textit{Parameter IDs} (PIDs) definidos pelo padrão OBD-II. 

Do ponto de vista de segurança, espera-se que a solução proposta incorpore mecanismos capazes de mitigar as vulnerabilidades identificadas em dispositivos comerciais. Entre os resultados esperados estão a implementação de autenticação entre o \textit{dongle} e o aplicativo móvel, proteção criptográfica da comunicação sem fio e controle de acesso às mensagens enviadas ao barramento CAN. Para isso, pretende-se avaliar o sistema por meio de testes de segurança inspirados nos métodos utilizados em trabalhos da literatura, verificando se as principais vulnerabilidades descritas são mitigadas pela arquitetura proposta.

Por fim, espera-se que este trabalho contribua para o avanço da segurança em dispositivos automotivos conectados, fornecendo uma referência de arquitetura e de boas práticas para o desenvolvimento de \textit{dongles} OBD-II e aplicações associadas com maior nível de proteção contra ataques.


\section{Organização do Trabalho}

A \autoref{cap:intro} explica como o OBD-II possibilita a coleta de dados do automóvel em tempo real, que por sua vez ensejam uma miríade de aplicações. Explicitam-se os objetivos de projeto e desenvolvimento de um \textit{dongle} e de um \textit{companion app} seguros.

\autoref{cap:conceitos} apresenta elementos conceituais relacionados ao trabalho, e analisa as principais referências da literatura utilizadas e também os conceitos de engenharia mais importantes aplicados no projeto. Preocupa-se também com a modelagem dos ataques.

No \autoref{cap:especificacao}, de especificação, são descritos de maneira técnicas e completa todos os requisitos funcionais e não funcionais que foram planejados para o \textit{dongle}  e para o \textit{companion app}  que serão desenvolvidos no trabalho.  Nele, também é introduzida a arquitetura do sistema.

Com a especificação do projeto melhor definida, o \autoref{cap:desenvolvimento} foca na sua execução prática, apresentando o cronograma detalhado de como o projeto será realizado dentro do tempo disponível para o trabalho.

Por fim, no \autoref{cap:consideracoes}, de considerações finais, o trabalho é encerrado, com destataque para as constribuições técnicas relevantes para o contexto da engenharia: reprodutibilidade, segurança \textit{open source} para estudos e melhorias dentro da comunidade de \textit{dongles} automotivos e contribuição acadêmica na implementação de um protocolo seguro para a interação com a rede do veículo.